% defer/updates.tex
% mainfile: ../perfbook.tex
% SPDX-License-Identifier: CC-BY-SA-3.0

\section{更新操作怎么办？}
\label{sec:defer:What About Updates?}
%
\epigraph{生活中唯一不变的就是变化。}
	 {Fran\c{c}ois de la Rochefoucauld}

本章所介绍的延迟处理技术最直接地适用于读多写少的场景，
这不禁让人要问：``那更新操作怎么办？''
毕竟，提升读者的性能和可扩展性固然很好，
但人们自然也希望写者能获得同样出色的性能和可扩展性。

我们已经见到了一种写者具有高性能和高可扩展性的情况，
即\cref{chp:Counting}中介绍的计数算法。
这些算法采用了部分分区的数据结构，使得更新操作可以在本地进行，
而开销较大的读取操作则必须对整个数据结构求和。
Silas Boyd-Wickizer 将这一概念推广为 OpLog，
并将其应用于 Linux 内核的路径名查找、虚拟内存反向映射和
\co{stat()} 系统调用~\cite{SilasBoydWickizerPhD}。

另一种方法称为``Disruptor''，专为处理大量输入数据流的应用而设计。
其方法是依靠单生产者-单消费者 FIFO 队列，
从而最大限度地减少同步需求~\cite{AdrianSutton2013LCA:Disruptor}。
对于 Java 应用，Disruptor 还具有减少垃圾回收器使用的优点。

当然，在可行的情况下，完全分区或``分片''的系统能提供出色的性能和可扩展性，
正如\cref{chp:Partitioning and Synchronization Design}中所述。

下一章将在几种不同类型的数据结构的背景下讨论更新操作。
